\chapter{Wechsel des Trägers}
\label{b13guide:ch:ambient}

\section{Schranken in einem kleineren Träger}

\begin{theorem}[Einschränken des umgebenden Trägers]
\label{b13guide:thm:restriction}
Seien $S\subseteq Q\subseteq P$, wobei $Q$ die von $P$
eingeschränkte Ordnung trägt. Dann gilt
\[
 \Lo_Q(S)=Q\cap\Lo_P(S),\qquad
 \Hi_Q(S)=Q\cap\Hi_P(S).
\]
Insbesondere ist Beschränktheit in $Q$ stärker als Beschränktheit
in $P$: Jede Schranke in $Q$ ist eine Schranke in $P$, aber eine
Schranke in $P$ muß nicht zu $Q$ gehören.
\end{theorem}

\begin{theorem}[Erhalten vorhandener Infima und Suprema]
\label{b13guide:thm:inheritance}
Seien $S\subseteq Q\subseteq P$. Existiert $s=\sup_P S$ und liegt
$s\in Q$, so existiert $\sup_Q S$, und $\sup_Q S=s$.
Existiert $m=\inf_P S$ und liegt $m\in Q$, so existiert $\inf_Q S$,
und $\inf_Q S=m$.
Die Existenz eines Supremums oder Infimums in $Q$ allein garantiert
weder dessen Existenz noch denselben Wert in $P$.
\end{theorem}

\begin{theorem}[Vergleich zwischen beiden Trägern]
\label{b13guide:thm:ambient-comparison}
Seien $S\subseteq Q\subseteq P$. Bei Existenz der jeweils beiden
Werte gilt
\[
 \sup_P S\leq\sup_Q S,\qquad \inf_Q S\leq\inf_P S.
\]
Ist $q=\sup_Q S$ vorhanden, so gilt
\[
 q=\sup_P S\quad\Longleftrightarrow\quad
       \forall u\in\Hi_P(S)\;(q\leq u).
\]
Ist $q=\inf_Q S$ vorhanden, so gilt
\[
 q=\inf_P S\quad\Longleftrightarrow\quad
       \forall \ell\in\Lo_P(S)\;(\ell\leq q).
\]
\end{theorem}


\begin{theorem}[Trägerwechsel bei ausreichend vielen Schranken]
\label{b13guide:thm:ambient-cofinal}
Seien $S\subseteq Q\subseteq P$. Gilt
\[
 \forall u\in\Hi_P(S)\;\exists v\in\Hi_Q(S)\;(v\leq u),
\]
so existiert $\sup_P S$ genau dann, wenn $\sup_Q S$ existiert;
im Existenzfall stimmen die Werte überein. Gilt die duale Bedingung
\[
 \forall\ell\in\Lo_P(S)\;\exists k\in\Lo_Q(S)\;(\ell\leq k),
\]
so existiert $\inf_P S$ genau dann, wenn $\inf_Q S$ existiert;
im Existenzfall stimmen die Werte überein.
Ist $q=\sup_Q S$ bereits vorhanden, so ist die erste Bedingung
auch notwendig für $q=\sup_P S$; für ein vorhandenes $q=\inf_Q S$
ist die zweite notwendig für $q=\inf_P S$.
Die Bedingungen betreffen die Schrankenmengen der betrachteten
Teilmenge $S$, nicht bloß die Lage von $Q$ in $P$.
\end{theorem}

